
\documentclass[12pt]{article}
\usepackage[a4paper, total={6.5in, 9.5in}]{geometry}
\usepackage{graphicx}			
\usepackage{amssymb}
\usepackage{amsmath}			
\usepackage{fancyhdr}
\usepackage{enumerate}       
\usepackage{dirtytalk}
\usepackage[english]{babel} 
\usepackage{amsthm}
\usepackage{xcolor}
\usepackage[shortlabels]{enumitem}
\setlength{\parindent}{0pt}
\setlength{\parskip}{0.8em}
\begin{document}
\renewcommand{\thesection}{Exercise \arabic{section}}
\parindent = 0pt


      {\bf Analysis 1 HW 1}  \hfill 


%==================

% exercise1
\section{}
 {\bf Definitions}


$f: S \rightarrow T \text{ is invertible} \Longleftrightarrow \exists ~g: T \rightarrow S \text{ with } fg = id_T \text{ and } gf = id_S$.


$f:S \rightarrow T$ is injective $\Longleftrightarrow f(x_1) = f(x_2) \implies x_1 = x_2 \ \ \forall \, x_1, \, x_2 \in S$.

$f:S \rightarrow T$ is surjective $\Longleftrightarrow \forall \, y \in T \ \ \exists \, x\in S  \  \text{with}  \  f(x) = y$.

\begin{enumerate}[(a)]
      \item
            {\bf P:} $f$ is invertible.

                  {\bf Q:} $f$ is injective and $f$ is surjective.

            To prove $P \Longleftrightarrow Q$, we need to prove $P \Longrightarrow Q$ and $\neg P \Longrightarrow \neg Q$.
            \begin{enumerate}[(i)]
                  \item $P \Longrightarrow Q$


                        \begin{proof}
                              From $P$, we have $g(f(x)) = x \ \ \forall \, x \in S$.
                              Consider $x_1, x_2 \in S$. Then,
                              \begin{align*}
                                    g(f(x_1)) & = x_1 \\
                                    g(f(x_2)) & = x_2
                              \end{align*}
                              Since $g$ is a function, $f(x_1) = f(x_2) \implies x_1 = x_2$.


                              $\therefore$ $f$ is injective.
                              \newline

                              Also from $P$, $f(g(y)) = y \ \ \forall \, y \in T$.
                              For this to be true, every element in $T$ must be
                              in the image of $f$.

                              $\therefore$ $f$ is surjective.

                        \end{proof}
                  \item $\neg P \Longrightarrow \neg Q$
                        \begin{proof}
                              \begin{align*}
                                    \neg P & \Longleftrightarrow \nexists \, g: T \rightarrow S \ \text{with} \ fg = id_T \text{ and } gf = id_S            \\
                                           & \Longleftrightarrow \forall \, g: T \rightarrow S \ fg \ne id_T \text{ or } gf \ne id_S                        \\
                                           & \Longleftrightarrow (\exists \, y \in T \ \text{with}\ y \notin f(S)) \text{ or }                              \\
                                           & \quad \quad \, \, \,(\exists \ x_1, x_2  \in S, \ x_1 \ne x_2, \ \text{with} \ f(x_1) = f(x_2) = y, \ y \in T) \\
                                           & \Longleftrightarrow (f \ \text{is not surjective}) \text{ or } (f \ \text{is not injective})                   \\
                                           & \Longleftrightarrow \neg Q
                              \end{align*}
                        \end{proof}
            \end{enumerate}
      \item
            Let $g$ and $h$ be inverses of $f$. Then,
            \begin{align*}
                   & f(g(x)) = x \ \ \forall \, x \in T \\
                   & f(h(x)) = x \ \ \forall \, x \in T \\ \\
                   & \Longrightarrow fg = fh
            \end{align*}
            If $f$ is invertible, from (a), $f$ is injective. $\therefore g = h$, i.e,
            the inverse of an invertible function is unique.
      \item

            \begin{enumerate}[(i)]
                  \item {\bf Left invertible functions:}
                        From (a)(i), $f: S \rightarrow T$ being left invertible is equivalent to $f$ being injective.
                        Let $g:T\rightarrow S$
                        be a left inverse of $f$.
                        While $f$ being injective guarantees that $g$ exists, $g$ will be unique if and only if $f$ is also surjective, as otherwise
                        $g(y),\ y \in T\setminus f(S)$ can be arbitrarily defined as any $x \in S$.


                  \item {\bf Right invertible functions:}
                        From (a)(i), $f: S \rightarrow T$ being right invertible implies $f$ is surjective. Let $g:T\rightarrow S$
                        be a right inverse of $f$. $g$ is unique if and only if $f$ is also injective, as otherwise $g(y),\, y\in T$ can be arbitrarily defined to
                        be any $x \in S$ such that $f(x) = y$.
                  \item {\bf Left and right invertible functions:}
                        If $f$ is left invertible and right invertible, it is injective and surjective,
                        which implies $f$ is invertible, which implies $f$ has a unique inverse, which must be
                        equal to its left and right inverses.
                  \item {\bf Left cancellable functions:} Let $h_1, h_2: T \rightarrow S$. For
                        $f:S \rightarrow T$
                        to be left cancellable,
                        \begin{align*}
                                                        & (fh_1 = fh_2) \implies (h_1 = h_2)                                  \\
                              \Longleftrightarrow \quad & (f(x_1) = f(x_2)) \implies (x_1 = x_2) \ \ \forall \,x_1, x_2 \in S
                        \end{align*}
                        which implies $f$ is injective. Evidently, the terms \say{injective}, \say{left invertible},
                        and \say{left cancellable} are equivalent.
                  \item {\bf Right cancellable functions:} Borrowing $h_1$, $h_2$, and $f$ from (iv), $f$
                        is right cancellable if
                        \begin{align*}
                                                        & (h_1f = h_2f) \implies (h_1 = h_2)                                                       \\
                              \Longleftrightarrow \quad & (h_1f(x) = h_2f(x)\ \ \forall \,x \in S) \implies (h_1(y) = h_2(y)\ \ \forall \,y \in T) \\
                              \Longleftrightarrow \quad & \forall \,y \in T \ \ \exists \,x \in X \ \text{with}\  f(x) = y
                        \end{align*}
                        which implies $f$ is surjective. Evidently, the terms \say{surjective}, \say{right invertible},
                        and \say{right cancellable} are equivalent.
                  \item {\bf Left and right cancellable functions:} If $f$ is left cancellable and right cancellable, it is injective and surjective,
                        which implies $f$ is invertible.
            \end{enumerate}
      \item The claim is no longer true if the condition of the domain and codomain being of
            finite cardinality is dropped. For example, $f:\mathbb{N} \rightarrow \mathbb{R}$
            defined as $f(x) = x$ is injective, and has a domain and a codomain both of infinite
            cardinality. However, one is countably infinite which the other is uncountably
            infinite, making it impossible for the function to be surjective.
\end{enumerate}

% exercise2
\section{}
 {\bf Definitions}\\
$S$ is the universal set.\\
$A$ is the set of all equivalence relations on S.\\
$B$ is the set of all partitions of S.\\
$[s]_a = \{k \in S \ | \ (k, s) \in a\}, s\in S, a \in A$ is the equivalence class of $s$ in $a$.\\
$f:A\rightarrow B$, $f(a) = \{[s]_a\ | \ s \in S\}$ maps equivalence relations to partitions.\\
$g:B\rightarrow A$, $g(b) = \{(n, m)\ | \exists \ e \in b \text{ with } n, m \in e\}$ maps partitions to equivalence classes.

      {\bf Claim}\\
$f$ and $g$ are inverses of each other.
\begin{enumerate}[(i)]
      \item $f(a), a\in A$ is a partition.
            \begin{proof}

                  \begin{enumerate}[(a)]
                        \item For $s_1, s_2 \in S, a \in A$ if $[s_1]_a \cap [s_2]_a \ne \phi$, there exists $x$ such that
                              $x \in [s_1]_a$ and $x \in [s_2]_a$. This implies $(x, s_1) \in a$ and $(x, s_2) \in a$. Since $a$
                              is symmetric, $(s_1, x) \in a$, which
                              implies $(s_1, s_2) \in a$. Thus, $[s_1]_a = [s_2]_a$. $\therefore$ the elements
                              of $f(a)$ are disjoint.

                        \item For every equivalence relation $a \in A$, $(s, s) \in a \ \ \forall\, s \in S$. Thus,
                              there exists $e \in f(a)$ such that $s \in e, \ \ \forall\, s \in S$. $\therefore$
                              $f(a)$ is exhaustive.
                  \end{enumerate}

                  $f(a)$ is disjoint and exhaustive, which implies that $f(a)$ is a partition of $S$.
            \end{proof}
      \item $g(b), b\in B$ is an equivalence relation.
            \begin{proof}

                  \begin{enumerate}[(a)]
                        \item $\exists e \in b$ such that $s\in e \ \ \forall s \in S$. Thus, $(s, s) \in g(b) \ \ \forall s \in S$.
                              $\therefore$ $g(b)$ is reflexive.

                        \item $s_1, s_2 \in e \Longleftrightarrow s_2, s_1 \in e$. Thus,
                              if $(s_1, s_2) \in g(b)$, $(s_2, s_1) \in g(b) \ \ \forall \, s_1, s_2 \in S$.
                              $\therefore$ $g(b)$ is symmetric.

                        \item $s_1, s_2 \in e_1$ and $s_2, s_3 \in e_2$ for some $e_1, e_2 \in b$,
                              $s_2 \in e_1 \cap e_2$, which implies that $e_1 = e_2$. Hence, If
                              $(s_1, s_2) \in g(b)$ and $(s_2, s_3) \in g(b)$, $(s_1, s_3) \in g(b)$.
                              $\therefore$ $g(b)$ is transitive.
                  \end{enumerate}
                  $g(b)$ is reflexive, symmetric, and transitive, which implies $g(b)$ is an equivalence
                  relation.
            \end{proof}
      \item $fg = id_B$
            \begin{proof}
                  Let $s_1, s_2 \in S$. Let $b \in B$, so $b$ is a partition of $S$. From the definition of $g$,
                  $(s_1, s_2) \in g(b)$ if and only if $s_1$ and $s_2$ belong to the same part of
                  $b$. From the definition of $f$, $\exists e \in f(g(b))$ such that $s_1, s_2 \in e$
                  if and only if $(s_1, s_2) \in g(b)$. Thus,
                  \begin{align*}
                                               & \exists e\in b \text{ such that } s_1, s_2 \in e       \\
                        \Longleftrightarrow \  & (s_1, s_2) \in g(b)                                    \\
                        \Longleftrightarrow \  & \exists e\in f(g(b)) \text{ such that } s_1, s_2 \in e
                  \end{align*}
                  $\therefore$ $f(g(b)) = b \ \ \forall b \in B$.
            \end{proof}
      \item $gf = id_A$
            \begin{proof}
                  Let $s_1, s_2 \in S$. Let $a \in A$. From the definition of $f$, $s_1$ and $s_2$ belong
                  to the same part of $f(a)$ if and only if $(s_1, s_2) \in a$. From the definition of $g$,
                  $(s_1, s_2) \in g(f(a))$ if and only if $s_1$ and $s_2$ belong to the same part of $f(a)$.
                  Thus,
                  \begin{align*}
                                               & (s_1, s_2) \in a                                    \\
                        \Longleftrightarrow \  & \exists e\in f(a) \text{ such that } s_1, s_2 \in e \\
                        \Longleftrightarrow \  & (s_1, s_2) \in g(f(a))
                  \end{align*}
                  $\therefore$ $g(f(a)) = a\ \ \forall a \in A$.
            \end{proof}
\end{enumerate}
From (iii) and (iv), $f$ is bijective. This implies that there exists a one to one
correspondence between the elements of $A$ and those of $B$. Thus, defining a partition of
$S$ is "the same as" defining a equivalence relation on $S$ in the sense that every partition
corresponds uniquely to an equivalence relation and vice versa.


% exercise3 
\section{}
 {\bf Definitions} \\
$f: S \rightarrow T$ \\
$f(A) = \{f(x) \ |\  x\in A \}$, $A \subset S$ \\
$f^{-1}(B) = \{x\in S \ |\  f(x) \in B\}$, $B \subset T$
\begin{enumerate}[(a)]
      \item $f(\bigcup_{i\in I} A_i) = \bigcup_{i\in I} f(A_i)$,
            $A_i \subset S$
            \begin{proof}
                  \begin{align*}
                                               & y \in f\Big(\bigcup_{i\in I} A_i\Big)                                        \\
                        \Longleftrightarrow \  & \exists A_i \text{ such that } \exists x \in A_i \text{ such that } f(x) = y \\
                        \Longleftrightarrow \  & \exists A_i \text{ such that } y \in f(A_i)                                  \\
                        \Longleftrightarrow \  & y \in \bigcup_{i\in I} f(A_i)
                  \end{align*}
            \end{proof}
      \item $f(\bigcap_{i\in I} A_i) \ne \bigcap_{i\in I} f(A_i)$
            $A_i \subset S$
            \begin{proof}
                  Consider an $x_i$ in every $A_i$, such that no $x_i$ is in $\bigcap_{i\in I} A_i$.
                  Let $f(x_i) = y \ \ \forall \, i \in I$, and $f(x) \ne y$ for all other $x \in S$. Clearly, $y \notin f(\bigcap_{i\in I} A_i)$.
                  However, $y \in f(A_i) \ \ \forall \, i \in I \implies y \in \bigcap_{i\in I} f(A_i).$
            \end{proof}
      \item $f(A^c) \ne f(A)^c$
            \begin{proof}
                  Consider $x_1, x_2 \in S$ such that $x_1 \in A$ and $x_2 \notin A$. Let $f(x_1)
                        = f(x_2) = y$. Now, $x_2 \in A^c \implies y \in f(A^c)$. However, $x_1 \in A \implies
                        y \in f(A) \implies y \notin f(A)^c$.
            \end{proof}
      \item \begin{enumerate}[(i)]

                  \item $f^{-1}(\bigcup_{i\in I} B_i) = \bigcup_{i\in I} f^{-1}(B_i)$,
                        $B_i \subset T$
                        \begin{proof}
                              \begin{align*}
                                                           & x \in f^{-1}\Big(\bigcup_{i\in I} B_i\Big)       \\
                                    \Longleftrightarrow \  & \exists B_i \text{ such that } f(x) \in B_i      \\
                                    \Longleftrightarrow \  & \exists B_i \text{ such that } x \in f^{-1}(B_i) \\
                                    \Longleftrightarrow \  & x \in \bigcup_{i\in I} f^{-1}(B_i)
                              \end{align*}
                        \end{proof}

                  \item $f^{-1}(\bigcap_{i\in I} B_i) = \bigcap_{i\in I} f^{-1}(B_i)$
                        $B_i \subset T$
                        \begin{proof}
                              \begin{align*}
                                                           & x \in f^{-1}\Big(\bigcap_{i\in I} B_i\Big) \\
                                    \Longleftrightarrow \  & f(x) \in B_i \ \ \forall \,i \in I         \\
                                    \Longleftrightarrow \  & x \in f^{-1}(B_i) \ \ \forall \, i \in I   \\
                                    \Longleftrightarrow \  & x \in \bigcap_{i\in I} f^{-1}(B_i)
                              \end{align*}
                        \end{proof}
                  \item $f^{-1}(B^c) = f^{-1}(B)^c$
                        \begin{proof}
                              \begin{align*}
                                                           & x \in f^{-1}(B^c)  \\
                                    \Longleftrightarrow \  & f(x) \in B^c       \\
                                    \Longleftrightarrow \  & f(x) \notin B      \\
                                    \Longleftrightarrow \  & x \notin f^{-1}(B) \\
                                    \Longleftrightarrow \  & x \in f^{-1}(B)^c
                              \end{align*}
                        \end{proof}
            \end{enumerate}
      \item
            \begin{enumerate}[(i)]
                  \item  $f^{-1}(f(A)) \ne A$.  \begin{proof}
                              Consider $x_1 \in A, \  x_2, \notin A,\  f(x_1) = f(x_2) = y$.
                              $y \in f(A)$, which implies $x_1, x_2 \in f^{-1}(f(A))$.
                        \end{proof}

                  \item$f(f^{-1}(B)) \ne B$.
                        \begin{proof}
                              Consider $y \in T$ such that $y$ is not in the image of $f$. Consider
                              $B \subset T$ such that $y \in B$. Since there does not exist $x$ in
                              $f^{-1}(B)$ such that $f(x) = y$, $y \notin f(f^{-1}(B))$
                        \end{proof}
            \end{enumerate}

      \item
            Parts (b), (c), and (e)(i) relied on $f$ not being injective to provide
            counterexamples. If $f$ is injective, it can be proved that $f$ preserves
            intersections and complements, like (d)(ii) and (d)(iii). In (e)(ii), $f(f^{-1}(B))= B$
            if $f$ is surjective.
\end{enumerate}
\end{document}